\scanpage{150}
% archive foliation: 149.
\[
g.\tau=\tau\circ\tau_\infty(g^{-1})
\qquad\text{où}\qquad
\begin{gathered}
\tau_\infty=\begin{pmatrix}-1&0\\0&1\end{pmatrix}\in\mathrm{Gl}(2,\mathbb Z)\\
\tau_\infty(g^{-1})=\begin{pmatrix}d&b\\c&a\end{pmatrix}
\quad\text{si}\quad
g=\begin{pmatrix}a&b\\c&d\end{pmatrix}
\end{gathered}
\]

\uline{Corollaire 1} Le groupe fondamental $\mathcal T_{1,1}$ de
$M_{11}$ relatif au revêtement universel précédent
est canoniquement isomorphe à $\mathrm{Sl}(2,\mathbb Z)$
(avec l'opération précédente \inserted{à gauche} de $\mathrm{Sl}(2,\mathbb Z)$ sur
ledit revêtement, déduite de l'opération
canonique \inserted{à} droite par le passage de $g$ à $g^{-1}$\unsure{\ldots}).

Les marteaux-pilons sur les $\mathcal T_{g,\unsure{\nu|n}}$
donnent donc:

\uline{Corollaire 2} Soit $(\unsure{X|E},o)$ une surface orientée
compacte de genre 1, $\mathcal T$ le groupe des
classes d'isotopie d'automorphismes de $(X,o)$,
$H=\pi_1(X,o)=H_1(X,\mathbb Z)$, \deleted{\unsure{on a donc}} \inserted{alors} l'homo-
morphisme canonique
\[
\mathcal T\longrightarrow\mathrm{Aut}_{\mathbb Z}^+(H)
\tag{16}
\]
est un isomorphisme.

\uline{Corollaire 3} Soient $\Pi$ un groupe à lacets
de type $(1,1)$, $H=\Pi_{\mathrm{ab}}$, $\mathcal T$ le groupe des
autom. ext. à lacets de $\Pi$ compatibles avec
l'orientation, alors l'hom canonique
\[
\mathcal T\longrightarrow\mathrm{Aut}_{\mathbb Z}^+(H)
\]
est un iso.
